We presented MCTS-PIE, a unified study of six planners for multi-objective pathfinding in interactive environments, evaluated end-to-end across four experiment families (E1--E4) on a uniform 30-seed protocol with full per-run Pareto archives, ablation-driven defaults, and Holm--Bonferroni-corrected statistics. Across the main comparison, anytime curves, ablations, and environment scaling, the A*+MCTS hybrid is the strongest method overall, but a tuned multi-objective two-phase MCTS closes most of the gap and becomes the more practical choice as grid size grows: at $50 \times 50$ the search-only MO two-phase variant matches the hybrid's HV within $30\%$ at roughly $1/75$ the runtime, and at the most demanding cell ($50 \times 50$, $10$ checkpoints) it is the only method whose runtime remains tractable while the A*-based methods take more than 17~minutes per run.

Three findings are particularly noteworthy. First, hypervolume-based tree selection is the single highest-impact algorithmic component (E3): it more than doubles HV over plain UCB, and once it is in place the per-node Pareto archive of MO-MCTS becomes essentially free relative to plain MCTS. Second, the two-phase decomposition is brittle without multi-objective phase-1 selection: a single-objective phase-1 collapses on long checkpoint chains because UCB has no incentive to maintain path diversity. Third, anytime behavior is sharply non-linear for two-phase methods -- below 4k simulations they are unusable, but they ``ignite'' at 4k and stay competitive thereafter.

Future work will extend this study to learned heuristics for the shift subproblem, larger maps with non-rectangular topology, and stochastic shift outcomes (where shifting weight succeeds only with some probability). The MCTS-PIE harness was designed for these extensions: any new algorithm only needs to expose a \texttt{(spec) -> front} interface to plug into the existing E1--E4 protocol.
